% Template for PLoS
% Version 3.5 March 2018
%
% % % % % % % % % % % % % % % % % % % % % %
%
% -- IMPORTANT NOTE
%
% This template contains comments intended
% to minimize problems and delays during our production
% process. Please follow the template instructions
% whenever possible.
%
% % % % % % % % % % % % % % % % % % % % % % %
%
% Once your paper is accepted for publication,
% PLEASE REMOVE ALL TRACKED CHANGES in this file
% and leave only the final text of your manuscript.
% PLOS recommends the use of latexdiff to track changes during review, as this will help to maintain a clean tex file.
% Visit https://www.ctan.org/pkg/latexdiff?lang=en for info or contact us at latex@plos.org.
%
%
% There are no restrictions on package use within the LaTeX files except that
% no packages listed in the template may be deleted.
%
% Please do not include colors or graphics in the text.
%
% The manuscript LaTeX source should be contained within a single file (do not use \input, \externaldocument, or similar commands).
%
% % % % % % % % % % % % % % % % % % % % % % %
%
% -- FIGURES AND TABLES
%
% Please include tables/figure captions directly after the paragraph where they are first cited in the text.
%
% DO NOT INCLUDE GRAPHICS IN YOUR MANUSCRIPT
% - Figures should be uploaded separately from your manuscript file.
% - Figures generated using LaTeX should be extracted and removed from the PDF before submission.
% - Figures containing multiple panels/subfigures must be combined into one image file before submission.
% For figure citations, please use "Fig" instead of "Figure".
% See http://journals.plos.org/plosone/s/figures for PLOS figure guidelines.
%
% Tables should be cell-based and may not contain:
% - spacing/line breaks within cells to alter layout or alignment
% - do not nest tabular environments (no tabular environments within tabular environments)
% - no graphics or colored text (cell background color/shading OK)
% See http://journals.plos.org/plosone/s/tables for table guidelines.
%
% For tables that exceed the width of the text column, use the adjustwidth environment as illustrated in the example table in text below.
%
% % % % % % % % % % % % % % % % % % % % % % % %
%
% -- EQUATIONS, MATH SYMBOLS, SUBSCRIPTS, AND SUPERSCRIPTS
%
% IMPORTANT
% Below are a few tips to help format your equations and other special characters according to our specifications. For more tips to help reduce the possibility of formatting errors during conversion, please see our LaTeX guidelines at http://journals.plos.org/plosone/s/latex
%
% For inline equations, please be sure to include all portions of an equation in the math environment.
%
% Do not include text that is not math in the math environment.
%
% Please add line breaks to long display equations when possible in order to fit size of the column.
%
% For inline equations, please do not include punctuation (commas, etc) within the math environment unless this is part of the equation.
%
% When adding superscript or subscripts outside of brackets/braces, please group using {}.
%
% Do not use \cal for caligraphic font.  Instead, use \mathcal{}
%
% % % % % % % % % % % % % % % % % % % % % % % %
%
% Please contact latex@plos.org with any questions.
%
% % % % % % % % % % % % % % % % % % % % % % % %

\documentclass[10pt,letterpaper]{article}
\usepackage[top=0.85in,left=2.75in,footskip=0.75in]{geometry}

% amsmath and amssymb packages, useful for mathematical formulas and symbols
\usepackage{amsmath,amssymb}

% Use adjustwidth environment to exceed column width (see example table in text)
\usepackage{changepage}


% textcomp package and marvosym package for additional characters
\usepackage{textcomp,marvosym}

% cite package, to clean up citations in the main text. Do not remove.
% \usepackage{cite}

% Use nameref to cite supporting information files (see Supporting Information section for more info)
\usepackage{nameref,hyperref}

% line numbers
\usepackage[right]{lineno}

% ligatures disabled
\usepackage{microtype}
\DisableLigatures[f]{encoding = *, family = * }

% color can be used to apply background shading to table cells only
\usepackage[table]{xcolor}

% array package and thick rules for tables
\usepackage{array}

% create "+" rule type for thick vertical lines
\newcolumntype{+}{!{\vrule width 2pt}}

% create \thickcline for thick horizontal lines of variable length
\newlength\savedwidth
\newcommand\thickcline[1]{%
  \noalign{\global\savedwidth\arrayrulewidth\global\arrayrulewidth 2pt}%
  \cline{#1}%
  \noalign{\vskip\arrayrulewidth}%
  \noalign{\global\arrayrulewidth\savedwidth}%
}

% \thickhline command for thick horizontal lines that span the table
\newcommand\thickhline{\noalign{\global\savedwidth\arrayrulewidth\global\arrayrulewidth 2pt}%
\hline
\noalign{\global\arrayrulewidth\savedwidth}}


% Remove comment for double spacing
%\usepackage{setspace}
%\doublespacing

% Text layout
\raggedright
\setlength{\parindent}{0.5cm}
\textwidth 5.25in
\textheight 8.75in

% Bold the 'Figure #' in the caption and separate it from the title/caption with a period
% Captions will be left justified
\usepackage[aboveskip=1pt,labelfont=bf,labelsep=period,justification=raggedright,singlelinecheck=off]{caption}
\renewcommand{\figurename}{Fig}

% Use the PLoS provided BiBTeX style
% \bibliographystyle{plos2015}

% Remove brackets from numbering in List of References
\makeatletter
\renewcommand{\@biblabel}[1]{\quad#1.}
\makeatother



% Header and Footer with logo
\usepackage{lastpage,fancyhdr,graphicx}
\usepackage{epstopdf}
%\pagestyle{myheadings}
\pagestyle{fancy}
\fancyhf{}
%\setlength{\headheight}{27.023pt}
%\lhead{\includegraphics[width=2.0in]{PLOS-submission.eps}}
\rfoot{\thepage/\pageref{LastPage}}
\renewcommand{\headrulewidth}{0pt}
\renewcommand{\footrule}{\hrule height 2pt \vspace{2mm}}
\fancyheadoffset[L]{2.25in}
\fancyfootoffset[L]{2.25in}
\lfoot{\today}

%% Include all macros below

\newcommand{\lorem}{{\bf LOREM}}
\newcommand{\ipsum}{{\bf IPSUM}}


% tightlist command for lists without linebreak
\providecommand{\tightlist}{%
  \setlength{\itemsep}{0pt}\setlength{\parskip}{0pt}}

% From pandoc table feature
\usepackage{longtable,booktabs,array}
\usepackage{calc} % for calculating minipage widths
% Correct order of tables after \paragraph or \subparagraph
\usepackage{etoolbox}
\makeatletter
\patchcmd\longtable{\par}{\if@noskipsec\mbox{}\fi\par}{}{}
\makeatother
% Allow footnotes in longtable head/foot
\IfFileExists{footnotehyper.sty}{\usepackage{footnotehyper}}{\usepackage{footnote}}
\makesavenoteenv{longtable}

% Pandoc citation processing
%From Pandoc 3.1.8
% definitions for citeproc citations
\NewDocumentCommand\citeproctext{}{}
\NewDocumentCommand\citeproc{mm}{%
  \begingroup\def\citeproctext{#2}\cite{#1}\endgroup}
\makeatletter
 % allow citations to break across lines
 \let\@cite@ofmt\@firstofone
 % avoid brackets around text for \cite:
 \def\@biblabel#1{}
 \def\@cite#1#2{{#1\if@tempswa , #2\fi}}
\makeatother
\newlength{\cslhangindent}
\setlength{\cslhangindent}{1.5em}
\newlength{\csllabelwidth}
\setlength{\csllabelwidth}{3em}
\newenvironment{CSLReferences}[2] % #1 hanging-indent, #2 entry-spacing
 {\begin{list}{}{%
  \setlength{\itemindent}{0pt}
  \setlength{\leftmargin}{0pt}
  \setlength{\parsep}{0pt}
  % turn on hanging indent if param 1 is 1
  \ifodd #1
   \setlength{\leftmargin}{\cslhangindent}
   \setlength{\itemindent}{-1\cslhangindent}
  \fi
  % set entry spacing
  \setlength{\itemsep}{#2\baselineskip}}}
 {\end{list}}
\usepackage{calc}
\newcommand{\CSLBlock}[1]{#1\hfill\break}
\newcommand{\CSLLeftMargin}[1]{\parbox[t]{\csllabelwidth}{#1}}
\newcommand{\CSLRightInline}[1]{\parbox[t]{\linewidth - \csllabelwidth}{#1}\break}
\newcommand{\CSLIndent}[1]{\hspace{\cslhangindent}#1}

\providecommand{\pandocbounded}[1]{#1}


\usepackage{forarray}
\usepackage{xstring}
\newcommand{\getIndex}[2]{
  \ForEach{,}{\IfEq{#1}{\thislevelitem}{\number\thislevelcount\ExitForEach}{}}{#2}
}

\setcounter{secnumdepth}{0}

\newcommand{\getAff}[1]{
  \getIndex{#1}{Independent researcher}
}

\begin{document}
\vspace*{0.2in}


% Title must be 250 characters or less.
\begin{flushleft}
{\Large
\textbf\newline{Frontier methods for digital soil mapping, evaluated
honestly: a reproducible benchmark on real WoSIS Cerrado
profiles} % Please use "sentence case" for title and headings (capitalize only the first word in a title (or heading), the first word in a subtitle (or subheading), and any proper nouns).
}
\newline
% Insert author names, affiliations and corresponding author email (do not include titles, positions, or degrees).
\\
Hugo Rodrigues\textsuperscript{\getAff{Independent
researcher}}\textsuperscript{*}\\
\bigskip
\textbf{\getAff{Independent researcher}}ORCID: 0000-0002-8070-8126\\
\bigskip
* Corresponding author: rodrigues.machado.hugo@gmail.com\\
\end{flushleft}
% Please keep the abstract below 300 words
\section*{Abstract}
Digital soil mapping (DSM) has moved well beyond regression trees, yet
the newer methods --- causal inference, physics-informed and generative
deep learning, graph and quantum kernels --- are rarely evaluated
against each other under a common, honest protocol, and their
predictions are rarely reported with calibrated uncertainty. We present
edaphos, an R framework that implements ten such method families behind
a single posterior class so that any of them can be scored with the same
diagnostics, and we use it to run a matched, spatially blocked benchmark
of six methods on 1095 real World Soil Information Service (WoSIS)
topsoil profiles over the Brazilian Cerrado. The result is deflationary
and, we argue, useful: a plain ordinary-least-squares model on a
directed-acyclic-graph (DAG)-selected covariate set is the most accurate
and best-calibrated method in the comparison, and none of the five more
elaborate pillars beats it on this task. We further show, on 3077
profiles, that backdoor adjustment yields robust associations with
explicit sensitivity bounds rather than identified causal effects, and
that instrumenting soil carbon with a trained foundation-model embedding
produces very strong but invalid instruments (first-stage F \(\approx\)
4630, Sargan p \(\approx\) 3e-13) --- an empirical confirmation that a
learned representation does not manufacture a valid instrument. We
conclude that the contribution of a framework like edaphos is not a new
best predictor but a reproducible, calibration-consistent apparatus for
asking which method helps, and for reporting honestly when none does.

% Please keep the Author Summary between 150 and 200 words
% Use first person. PLOS ONE authors please skip this step.
% Author Summary not valid for PLOS ONE submissions.

\linenumbers

% Use "Eq" instead of "Equation" for equation citations.
\section{Introduction}\label{introduction}

Digital soil mapping (DSM) predicts soil properties at unvisited
locations from spatially exhaustive covariates that proxy the
soil-forming factors {[}1{]}. The field's workhorse --- the random
forest and its quantile variants --- is accurate, fast and well
understood, and a large recent literature proposes to improve on it with
methods borrowed from causal inference, deep learning, and quantum
machine learning. Two problems recur. First, these methods are seldom
compared under a common protocol: differences in cross-validation
scheme, ensemble budget, target transform and evaluation metric make
published gains hard to attribute to the method rather than the harness.
Second, point accuracy is reported far more often than calibrated
uncertainty, even though a soil map is used to make decisions under that
uncertainty.

A third, more conceptual problem concerns causality. It is common in DSM
to read a covariate's importance in a predictive model as evidence that
it \emph{causes} the soil property {[}2{]}. Zhang and Wadoux {[}3{]}
argue that causal inference in DSM is possible only under conditions ---
an explicit causal model, causal sufficiency, and faithfulness --- that
observational soil data rarely satisfy, and that the generative,
process-informed view of causation is the one most compatible with those
conditions.

This article makes three contributions, all reproducible from data
bundled with the software: (i) a single uncertainty-calibration
interface that lets ten method families be scored with the same
diagnostics; (ii) a matched, spatially blocked, uncertainty-aware
benchmark of six methods on real WoSIS Cerrado profiles, whose honest
conclusion is that a DAG-informed OLS is the strongest method here; and
(iii) a causal analysis that reports backdoor-adjusted associations with
Cinelli--Hazlett robustness values, demonstrates confounding directly,
and reports a negative instrumental-variable result. We deliberately
foreground the negative findings: a framework that makes it easy to run
a dozen methods is only useful if it is equally easy to see when they do
not help.

\section{Materials and methods}\label{materials-and-methods}

\subsection{Data}\label{data}

The benchmark uses 1095 WoSIS {[}4{]} topsoil (0--30 cm) organic-carbon
profiles over the Cerrado, retrieved from the ISRIC Web Feature Service,
with covariates extracted at each profile from WorldClim 2.1 {[}5{]},
SoilGrids {[}6{]} and Copernicus fractional tree cover. The causal
analysis uses a larger draw of 3077 profiles over the same region and
covariate stack. Covariates are all open-licensed and fetched from
public endpoints, so the analysis is regenerable from profile
coordinates alone.

\subsection{A unified posterior and calibration
interface}\label{a-unified-posterior-and-calibration-interface}

Every predictive method in edaphos returns, or can be coerced to, an
\texttt{edaphos\_posterior}: a small object carrying a predictive sample
(or mean and standard deviation) for each target. A single routine then
scores any of them with the same measures --- the continuous ranked
probability score (CRPS), 90\% prediction-interval coverage (PICP) and
mean width (MPIW), and a reliability curve. This makes a like-for-like
comparison across otherwise incompatible methods possible: a random
forest's quantiles, a Gibbs sampler's draws and a diffusion model's
samples are all reduced to the same currency.

\subsection{Methods compared}\label{methods-compared}

The benchmark compares one baseline and five frontier methods:
\textbf{P1}, ordinary least squares on the covariates in the backdoor
adjustment set of an expert pedogenetic DAG (the baseline); \textbf{P4},
a quantile random forest on the principal components of a MoCo-v2
contrastive encoder trained on raster patches; \textbf{P5}, a quantile
random forest on the raw covariates; \textbf{P6}, quantum kernel ridge
regression using a ZZ-feature-map state-vector simulator; \textbf{P7}, a
Bayesian hierarchical spatial model fit by a Gibbs sampler; and
\textbf{P10}, a graph attention network over a
\emph{k}-nearest-neighbour co-location graph of profiles.

\subsection{Evaluation and causal
protocol}\label{evaluation-and-causal-protocol}

All methods are evaluated on the same 1095 spatially blocked
cross-validation folds, on the log1p scale of organic carbon, and scored
against the achievable null --- each fold's training-fold mean, not the
global mean. Ensemble size is reported rather than forced equal. For the
causal analysis we build an expert Cerrado DAG, compute the minimal
backdoor adjustment set per exposure, estimate the adjusted association
by OLS with a spatial block bootstrap (500 resamples over 20 clusters),
and pair every estimate with a Cinelli--Hazlett robustness value (RV)
{[}7{]}: the fraction of residual variance an unobserved confounder
would have to explain in both exposure and outcome to drive the estimate
to zero. We also attempt instrumental-variable estimation, using the
principal components of a foundation-model embedding as instruments for
precipitation, and test their validity with the Sargan
over-identification test.

\section{Results}\label{results}

\subsection{A DAG-informed OLS is the strongest
method}\label{a-dag-informed-ols-is-the-strongest-method}

\begin{longtable}[]{@{}
  >{\raggedright\arraybackslash}p{(\linewidth - 12\tabcolsep) * \real{0.4400}}
  >{\raggedleft\arraybackslash}p{(\linewidth - 12\tabcolsep) * \real{0.0800}}
  >{\raggedleft\arraybackslash}p{(\linewidth - 12\tabcolsep) * \real{0.0933}}
  >{\raggedleft\arraybackslash}p{(\linewidth - 12\tabcolsep) * \real{0.0933}}
  >{\raggedleft\arraybackslash}p{(\linewidth - 12\tabcolsep) * \real{0.0933}}
  >{\raggedleft\arraybackslash}p{(\linewidth - 12\tabcolsep) * \real{0.0800}}
  >{\raggedleft\arraybackslash}p{(\linewidth - 12\tabcolsep) * \real{0.1200}}@{}}
\caption{Matched cross-pillar benchmark on 1095 WoSIS Cerrado profiles
(1095 spatially blocked folds, log1p scale). RMSE and CRPS lower is
better; R2 and PICP near 0.90 better.}\tabularnewline
\toprule\noalign{}
\begin{minipage}[b]{\linewidth}\raggedright
Method
\end{minipage} & \begin{minipage}[b]{\linewidth}\raggedleft
RMSE
\end{minipage} & \begin{minipage}[b]{\linewidth}\raggedleft
R2
\end{minipage} & \begin{minipage}[b]{\linewidth}\raggedleft
PICP90
\end{minipage} & \begin{minipage}[b]{\linewidth}\raggedleft
MPIW90
\end{minipage} & \begin{minipage}[b]{\linewidth}\raggedleft
CRPS
\end{minipage} & \begin{minipage}[b]{\linewidth}\raggedleft
Time (s)
\end{minipage} \\
\midrule\noalign{}
\endfirsthead
\toprule\noalign{}
\begin{minipage}[b]{\linewidth}\raggedright
Method
\end{minipage} & \begin{minipage}[b]{\linewidth}\raggedleft
RMSE
\end{minipage} & \begin{minipage}[b]{\linewidth}\raggedleft
R2
\end{minipage} & \begin{minipage}[b]{\linewidth}\raggedleft
PICP90
\end{minipage} & \begin{minipage}[b]{\linewidth}\raggedleft
MPIW90
\end{minipage} & \begin{minipage}[b]{\linewidth}\raggedleft
CRPS
\end{minipage} & \begin{minipage}[b]{\linewidth}\raggedleft
Time (s)
\end{minipage} \\
\midrule\noalign{}
\endhead
\bottomrule\noalign{}
\endlastfoot
P1 OLS (DAG-selected covariates) & 0.583 & 0.067 & 0.878 & 1.844 & 0.323
& 2.3 \\
P5 QRF & 0.589 & 0.051 & 0.872 & 1.830 & 0.321 & 1.6 \\
P7 BHS & 0.597 & 0.041 & 0.845 & 1.692 & 0.334 & 111.1 \\
P4 Foundation embeddings + QRF & 0.613 & -0.003 & 0.884 & 2.143 & 0.336
& 80.6 \\
P10 GAT (trained attention) & 0.621 & -0.080 & 0.745 & 1.442 & 0.367 &
36.1 \\
P6 Quantum KRR & 0.663 & -0.184 & 0.476 & 0.863 & 0.417 & 715.4 \\
\end{longtable}

Ordinary least squares on the DAG-selected covariates has the lowest
RMSE (0.583) and the highest coefficient of determination (0.067), and
its calibration (PICP 0.878) is among the best (Table 1). The random
forest is a close second. The three deep or exotic methods ---
foundation embeddings, graph attention, and the quantum kernel --- sit
progressively lower, and all three fall below the achievable null
(\(R^2 < 0\)): on these spatially blocked folds they explain less than
predicting the training mean. The quantum kernel is both least accurate
and, at 715 s, several hundred times slower than the baseline (Fig 1).
The spread among the top methods is small relative to the fold-to-fold
standard deviation (\(\approx\) 0.06 RMSE), so we do not over-read the
ordering beyond its headline: elaboration did not pay here.

\begin{figure}
\centering
\pandocbounded{\includegraphics[keepaspectratio]{manuscript_files/manuscript_files/figure-latex/bench-plot-1.pdf}}
\caption{Accuracy (RMSE, lower better) against variance explained (R2 vs
the training-fold mean). The DAG-informed OLS (filled) is the most
accurate; every more elaborate method sits down and to the left, and
three fall below the achievable null.}
\end{figure}

When tuned under nested cross-validation (1095 profiles; outer folds
never used for tuning), the graph attention network reaches a pooled
Nash--Sutcliffe efficiency of 0.184 against the tuned forest's 0.165 ---
an edge whose 95\% block-bootstrap confidence interval, {[}-0.021,
0.057{]}, spans zero (P(GAT ahead) = 0.844). The attention mechanism is
thus competitive with, not superior to, a random forest.

\subsection{Causal inference: robust associations, and why
identification is
hard}\label{causal-inference-robust-associations-and-why-identification-is-hard}

\begin{longtable}[]{@{}
  >{\raggedright\arraybackslash}p{(\linewidth - 12\tabcolsep) * \real{0.2209}}
  >{\raggedright\arraybackslash}p{(\linewidth - 12\tabcolsep) * \real{0.2674}}
  >{\raggedleft\arraybackslash}p{(\linewidth - 12\tabcolsep) * \real{0.0814}}
  >{\raggedleft\arraybackslash}p{(\linewidth - 12\tabcolsep) * \real{0.1047}}
  >{\centering\arraybackslash}p{(\linewidth - 12\tabcolsep) * \real{0.1977}}
  >{\centering\arraybackslash}p{(\linewidth - 12\tabcolsep) * \real{0.0698}}
  >{\raggedleft\arraybackslash}p{(\linewidth - 12\tabcolsep) * \real{0.0581}}@{}}
\caption{Backdoor-adjusted associations on 3077 WoSIS Cerrado profiles
(g/kg topsoil carbon per unit exposure). RV: Cinelli-Hazlett robustness
value.}\tabularnewline
\toprule\noalign{}
\begin{minipage}[b]{\linewidth}\raggedright
\end{minipage} & \begin{minipage}[b]{\linewidth}\raggedright
Exposure
\end{minipage} & \begin{minipage}[b]{\linewidth}\raggedleft
Naive
\end{minipage} & \begin{minipage}[b]{\linewidth}\raggedleft
Adjusted
\end{minipage} & \begin{minipage}[b]{\linewidth}\centering
95\% CI
\end{minipage} & \begin{minipage}[b]{\linewidth}\centering
Sig.
\end{minipage} & \begin{minipage}[b]{\linewidth}\raggedleft
RV
\end{minipage} \\
\midrule\noalign{}
\endfirsthead
\toprule\noalign{}
\begin{minipage}[b]{\linewidth}\raggedright
\end{minipage} & \begin{minipage}[b]{\linewidth}\raggedright
Exposure
\end{minipage} & \begin{minipage}[b]{\linewidth}\raggedleft
Naive
\end{minipage} & \begin{minipage}[b]{\linewidth}\raggedleft
Adjusted
\end{minipage} & \begin{minipage}[b]{\linewidth}\centering
95\% CI
\end{minipage} & \begin{minipage}[b]{\linewidth}\centering
Sig.
\end{minipage} & \begin{minipage}[b]{\linewidth}\raggedleft
RV
\end{minipage} \\
\midrule\noalign{}
\endhead
\bottomrule\noalign{}
\endlastfoot
wc\_bio\_12 & Precipitation (per mm) & 0.008 & 0.010 & {[}0.003,
0.019{]} & yes & 0.16 \\
wc\_bio\_01 & Temperature (per degC) & -0.467 & -0.824 & {[}-1.367,
0.222{]} & no & 0.08 \\
wc\_landcover\_trees & Tree cover (per 0-1) & 2.864 & 0.347 & {[}-3.314,
2.940{]} & no & 0.01 \\
soilgrids\_clay & Clay (per \%) & 0.446 & 0.448 & {[}0.292, 0.641{]} &
yes & 0.18 \\
\end{longtable}

Two associations are robust and mechanistically sensible (Table 2):
precipitation and clay content both raise topsoil carbon, are
distinguishable from zero, and would each need an unobserved confounder
explaining about 16--18\% of residual variance in both exposure and
outcome to overturn. The temperature association carries the expected
negative sign but is not distinguishable from zero. The tree-cover row
is instructive: its naive association (2.86) all but vanishes after
adjusting for climate (0.35, interval spanning zero, RV 0.006) --- the
raw signal was largely shared climatic ancestry. This is exactly the
confounding that Zhang and Wadoux {[}3{]} warn is misread as causal,
shown directly on real data.

Backdoor adjustment identifies an effect only under causal sufficiency.
Instrumental variables promise identification under a different
assumption --- that an instrument affects the outcome only through the
exposure. We trained a real MoCo-v2 foundation encoder (7 covariate
channels, 12,000 tiles, 3,000 epochs) and used the principal components
of its per-profile landscape embedding as instruments for precipitation.
The instruments are extremely relevant --- a first-stage \(F\) of 4630,
far above the usual threshold of 10 --- yet they decisively fail the
Sargan over-identification test (\(\chi^2\)(4) = 84.4, \(p \approx\)
0e+00). A learned landscape representation encodes climate, terrain,
soil and vegetation together, so it reaches carbon through many paths
and violates the exclusion restriction. Even a powerful learned
representation does not manufacture a valid instrument; we report the
failed test, not a headline effect.

\section{Discussion}\label{discussion}

The benchmark's conclusion --- that a DAG-informed OLS is the strongest
method, and that three of five frontier pillars fall below the
achievable null --- deserves care in either direction. It does not show
that deep learning cannot work in DSM; it shows that on about 1000 point
profiles with a handful of coarse covariates, the sample and signal are
too small for high-capacity models to express an advantage, and that a
linear model on well-chosen covariates is hard to beat. This is a common
finding for tabular problems at this scale, and it is precisely the kind
of result that a matched, uncertainty-aware protocol surfaces and an
accuracy-only, single-method report hides. The calibration columns make
the same point from another angle: the quantum kernel's narrow intervals
buy the worst coverage, an overconfidence a point-accuracy table would
never reveal.

The causal results carry the second message. Backdoor adjustment on an
explicit DAG is a legitimate and useful tool --- it recovered a robust
precipitation and clay signal and exposed the tree-cover association as
mostly confounding --- but the honest object it returns is an adjusted
association with a sensitivity bound, not an identified causal effect.
The negative IV result sharpens this: the most tempting modern shortcut
to identification, a learned representation as instrument, is the wrong
tool, because the property that makes an embedding a good predictor (it
captures everything) is the property that makes it an invalid instrument
(it violates exclusion).

The limitations are inherent to the design. A single biome, about
1000--3000 profiles, and coarse covariates are a hard setting for
high-capacity methods; a larger, denser survey with fine-resolution
covariates might let the foundation encoder or graph network express
structure they cannot here, and the framework is built to make that
re-test a one-line change of data. The benchmark scores prediction, not
the process-informed generative modelling that Zhang and Wadoux {[}3{]}
identify as the most defensible route to causal claims.

\section{Conclusion}\label{conclusion}

We have presented a reproducible framework that reduces ten families of
DSM method to a common posterior-and-calibration currency, and used it
to run a matched, honest benchmark whose lesson is deflationary: on real
WoSIS Cerrado profiles a DAG-informed ordinary-least-squares model is
the most accurate and best-calibrated method, and no more elaborate
pillar beats it. The value we claim is not a new best predictor but the
apparatus --- and the discipline --- for finding out which method helps,
reporting calibrated uncertainty, and stating plainly when the answer is
``none of the new ones, here.'' The same apparatus delivered a robust
causal association analysis with explicit sensitivity bounds and a clean
negative instrumental-variable result, operationalising the caution that
causal identification from observational soil data is hard.

\section*{Data and code availability}\label{dataavail}
\addcontentsline{toc}{section}{Data and code availability}

edaphos and all bundled result data are available at
https://github.com/HugoMachadoRodrigues/edaphos. Soil profiles are from
WoSIS (ISRIC, CC-BY-4.0); covariates from WorldClim 2.1, SoilGrids and
Copernicus. The retrieval and analysis pipeline is in
\texttt{data-raw/}, and every figure in this manuscript is reproduced by
the \texttt{edaphos-paper} package vignette from the shipped data.

\section*{References}\label{references}
\addcontentsline{toc}{section}{References}

\protect\phantomsection\label{refs}
\begin{CSLReferences}{0}{1}
\bibitem[\citeproctext]{ref-McBratney2003}
\CSLLeftMargin{1. }%
\CSLRightInline{McBratney AB, Mendonça Santos ML, Minasny B. On digital
soil mapping. Geoderma. 2003;117: 3--52.
doi:\href{https://doi.org/10.1016/S0016-7061(03)00223-4}{10.1016/S0016-7061(03)00223-4}}

\bibitem[\citeproctext]{ref-Wadoux2020}
\CSLLeftMargin{2. }%
\CSLRightInline{Wadoux AMJ-C, Minasny B, McBratney AB. Machine learning
for digital soil mapping: Applications, challenges and suggested
solutions. Earth-Science Reviews. 2020;210: 103359.
doi:\href{https://doi.org/10.1016/j.earscirev.2020.103359}{10.1016/j.earscirev.2020.103359}}

\bibitem[\citeproctext]{ref-Zhang2026causal}
\CSLLeftMargin{3. }%
\CSLRightInline{Zhang L, Wadoux AMJ-C. Can digital soil mapping be
causal? European Journal of Soil Science. 2026;77: e70284.
doi:\href{https://doi.org/10.1111/ejss.70284}{10.1111/ejss.70284}}

\bibitem[\citeproctext]{ref-Batjes2024wosis}
\CSLLeftMargin{4. }%
\CSLRightInline{Batjes NH, Calisto L, Sousa LM de. Providing
quality-assessed and standardised soil data to support global mapping
and modelling ({WoSIS} snapshot 2023). Earth System Science Data.
2024;16: 4735--4765.
doi:\href{https://doi.org/10.5194/essd-16-4735-2024}{10.5194/essd-16-4735-2024}}

\bibitem[\citeproctext]{ref-Fick2017worldclim}
\CSLLeftMargin{5. }%
\CSLRightInline{Fick SE, Hijmans RJ. {WorldClim 2}: New 1-km spatial
resolution climate surfaces for global land areas. International Journal
of Climatology. 2017;37: 4302--4315.
doi:\href{https://doi.org/10.1002/joc.5086}{10.1002/joc.5086}}

\bibitem[\citeproctext]{ref-Hengl2017soilgrids}
\CSLLeftMargin{6. }%
\CSLRightInline{Hengl T, Mendes de Jesus J, Heuvelink GBM, Ruiperez
Gonzalez M, Kilibarda M, Blagotić A, et al. {SoilGrids250m}: Global
gridded soil information based on machine learning. PLOS ONE. 2017;12:
e0169748.
doi:\href{https://doi.org/10.1371/journal.pone.0169748}{10.1371/journal.pone.0169748}}

\bibitem[\citeproctext]{ref-CinelliHazlett2020}
\CSLLeftMargin{7. }%
\CSLRightInline{Cinelli C, Hazlett C. Making sense of sensitivity:
Extending omitted variable bias. Journal of the Royal Statistical
Society: Series B. 2020;82: 39--67.
doi:\href{https://doi.org/10.1111/rssb.12348}{10.1111/rssb.12348}}

\end{CSLReferences}

\nolinenumbers



\end{document}
